\honest{Where Philosophy Meets Science}{Eliezer Yudkowsky}{2008}

I look back on the ``dark ages'' of quantum physics, not to blame anyone but to learn a lesson. The most basic error, I say, is that the early scientists ``forgot that they themselves were made out of particles.'' Most of them knew it in theory, but they did not notice that a sensor detecting an electron, or even knowing about the electron's path, is just more ``particles in different places.'' So they brought in consciousness when they did not need it. I name none of them. \nb{The views that gave the observer's mind a role are usually traced to John von Neumann (1932), Fritz London and Edmond Bauer (1939), and Eugene Wigner (1961). Von Neumann allowed that the measuring instrument could be described by ordinary quantum evolution; he ended the chain of instruments with an ``abstract ego.'' Bohr rejected any role for the observer's mind.}

Our ancestors predicted other people by what they knew, so ``knowing'' feels simpler than brain states, as Thor once felt simpler than Maxwell's equations. Nobody said in physical terms when an observation occurs; the scientists ``just knew,'' so the theory could not be run as a computer program. \nb{The physicist John Bell made the same complaint in 1990: ``What exactly qualifies some physical systems to play the role of `measurer'?''}

The insight they needed looks philosophical. Yet ``it wasn't professional philosophers who swooped in and solved the problem and cleared up the mystery and made everything normal again. It was, well, physicists.'' \nb{Hugh Everett was a physics graduate student when he proposed his theory. But a month later, in ``Many Worlds, One Best Guess,'' Yudkowsky wrote that ``There is no known reason for the Born probabilities,'' and later work on that question includes philosophers of physics such as David Wallace, Hilary Greaves and Simon Saunders. Physicists still dispute whether many-worlds is right.}

Even when one philosopher gets it right in advance, I say, it is ``usually science'' that tells us which one. \nb{In ``The Dilemma: Science or Bayes?'', a month later, Yudkowsky rested many-worlds on simplicity, wrote that ``the evidence itself'' is not in dispute, and said that physicists who reject it ``aren't disobeying the rules of Science.'' So in the post's own case, no experiment has yet said which view won.}

It was once said that philosophy is what we have not yet turned into science. \nb{Bertrand Russell wrote a version of this in 1912.} I suggest that an insight looks philosophical when it is not yet taught in classes or reduced to calculation. The scientific method itself was once a philosophical notion, proved by the power of science and not by the agreement of philosophers.

So philosophy matters, but it is done well only ``from within'' a science. Professional philosophers should learn all the science they can, expect experiment to settle their arguments, and run experiments themselves. ``That, I think, is the lesson of history.''
